\documentclass[9pt]{article}
\usepackage{fontspec}
\usepackage[paperwidth=13.333in,paperheight=7.5in,margin=0.34in]{geometry}
\usepackage{amsmath}
\usepackage{xcolor}
\usepackage{tikz}
\usepackage{enumitem}
\usetikzlibrary{arrows.meta,positioning,calc}

\setmainfont{Apple SD Gothic Neo}
\setsansfont{Apple SD Gothic Neo}
\setmonofont{Menlo}
\pagestyle{empty}
\setlength{\parindent}{0pt}
\setlength{\parskip}{0pt}
\linespread{0.98}

\definecolor{TitleBlue}{HTML}{1D4ED8}
\definecolor{AccentOrange}{HTML}{EA580C}
\definecolor{SoftBlue}{HTML}{EFF6FF}
\definecolor{SoftGreen}{HTML}{ECFDF5}
\definecolor{SoftGray}{HTML}{F8FAFC}
\definecolor{SoftYellow}{HTML}{FEF3C7}
\definecolor{SoftWhite}{HTML}{FFFFFF}
\definecolor{LineGray}{HTML}{475569}
\definecolor{SoftBorder}{HTML}{CBD5E1}

\newcommand{\panel}[4]{
  \begin{tikzpicture}
    \node[
      draw=#2,
      fill=#3,
      rounded corners=7pt,
      line width=0.95pt,
      inner sep=7pt,
      text width=#4,
      align=left
    ] {#1};
  \end{tikzpicture}
}

\begin{document}
\sffamily

{\fontsize{19}{22}\selectfont\bfseries\color{TitleBlue} DVL의 \(\omega \times r\) 회전 성분 보정}

\vspace{0.03in}
{\small MuJoCo DVL raw 속도에서 \textbf{센서 오프셋 때문에 회전만으로 생기는 추가 속도}를 빼서 body 중심 기준 속도에 가깝게 만든다.}

\vspace{0.04in}

\begin{minipage}[t]{0.43\linewidth}
  \panel{
    {\large\bfseries\color{TitleBlue} 1. 핵심 개념}

    \vspace{0.04in}
    {\small
    \[
      \mathbf{v}_{sensor} = \mathbf{v}_{body} + \boldsymbol{\omega}\times\mathbf{r},
      \quad
      \mathbf{v}_{body} \approx \mathbf{v}_{dvl,raw} - \boldsymbol{\omega}\times\mathbf{r}
    \]

    \begin{itemize}[leftmargin=1.25em,itemsep=0.18em,topsep=0.25em]
      \item \(\boldsymbol{\omega}\): 기체 각속도
      \item \(\mathbf{r}\): 기체 중심에서 DVL까지의 위치 벡터
      \item \(\boldsymbol{\omega}\times\mathbf{r}\): 회전만으로 센서 위치에서 생기는 추가 속도
    \end{itemize}
    }
  }{TitleBlue}{SoftBlue}{0.96\linewidth}

  \vspace{0.04in}

  \panel{
    {\large\bfseries\color{AccentOrange} 2. 왜 빼야 하나}

    \vspace{0.04in}
    {\small
    \begin{itemize}[leftmargin=1.25em,itemsep=0.28em,topsep=0.25em]
      \item DVL은 기체 중심이 아니라 \texttt{tank\_current\_scene.xml}의 \texttt{dvl\_site}처럼 \textbf{오프셋된 위치}에 있다.
      \item 기체가 제자리에서 yaw, pitch, roll만 해도 센서 위치는 원운동한다.
      \item 이 성분을 그대로 두면 \textbf{회전을 병진 속도로 오해}해서 DVL odometry가 틀어진다.
    \end{itemize}
    }
  }{AccentOrange}{SoftYellow}{0.96\linewidth}

  \vspace{0.04in}

  \panel{
    {\large\bfseries\color{green!40!black} 3. 실제 구현}

    \vspace{0.03in}
    {\small
    \begin{itemize}[leftmargin=1.25em,itemsep=0.22em,topsep=0.25em]
      \item MuJoCo scene에서 \texttt{velocimeter dvl\_vel\_body}와 \texttt{rangefinder dvl\_altitude}를 읽는다.
      \item bridge가 DVL 센서축을 body 축으로 돌린다.
      \item 그 다음 \texttt{vel\_body = vel\_body - np.cross(gyro\_bmj, r\_body)}로 회전 성분을 뺀다.
      \item 보정값을 \texttt{/dvl/velocity}, \texttt{/dvl/odometry}, \texttt{/dvl/data}로 publish 한다.
    \end{itemize}
    }
  }{green!45!black}{SoftGreen}{0.96\linewidth}
\end{minipage}
\hfill
\begin{minipage}[t]{0.53\linewidth}
  \panel{
    {\large\bfseries\color{LineGray} 직관 그림}

    \vspace{0.01in}
    \begin{center}
    \begin{tikzpicture}[x=0.92in,y=0.92in]
      \draw[rounded corners=9pt, fill=SoftGray, draw=SoftBorder, line width=0.95pt] (0.0,0.0) rectangle (5.25,2.85);

      \coordinate (C) at (1.80,1.55);
      \coordinate (S) at (3.45,1.00);

      \filldraw[fill=TitleBlue, draw=TitleBlue] (C) circle (0.055);
      \node[anchor=south west, text=TitleBlue, font=\bfseries] at ($(C)+(0.08,0.00)$) {Body origin};

      \filldraw[fill=AccentOrange, draw=AccentOrange] (S) circle (0.055);
      \node[anchor=west, text=AccentOrange, font=\bfseries] at ($(S)+(0.10,-0.02)$) {DVL sensor};

      \draw[-{Latex[length=3.4mm]}, line width=1.20pt, color=LineGray] (C) -- (S)
        node[midway, above, sloped, fill=white, inner sep=1.5pt] {\(\mathbf{r}\)};

      \draw[-{Latex[length=3.5mm]}, line width=1.35pt, color=TitleBlue]
        ($(C)+(135:0.45)$) arc[start angle=135,end angle=20,radius=0.45];
      \node[text=TitleBlue, font=\bfseries] at ($(C)+(0.56,0.68)$) {\(\boldsymbol{\omega}\)};

      \draw[-{Latex[length=3.6mm]}, line width=1.45pt, color=AccentOrange] (S) -- ++(0.92,0.48)
        node[midway, above, fill=white, inner sep=1.5pt] {\(\boldsymbol{\omega}\times\mathbf{r}\)};

      \draw[-{Latex[length=3.6mm]}, line width=1.55pt, color=green!45!black] (S) -- ++(0.08,1.02)
        node[midway, right, fill=white, inner sep=1.5pt] {\(\mathbf{v}_{dvl,raw}\)};

      \draw[-{Latex[length=3.6mm]}, line width=1.55pt, color=purple!70!black] (C) -- ++(0.18,0.92)
        node[midway, left, fill=white, inner sep=1.5pt] {\(\mathbf{v}_{body}\)};

      \node[
        draw=SoftBorder,
        rounded corners=8pt,
        fill=SoftWhite,
        align=left,
        text width=2.20in,
        inner sep=6pt
      ] at (3.58,2.30) {
        \small
        \textbf{해석}\\
        기체가 제자리에서 도는 동안에도 DVL 위치에서는 속도가 생긴다.\\
        따라서 raw DVL 속도에서 \(\omega \times r\)를 빼야 body 중심 속도에 가까워진다.
      };
    \end{tikzpicture}
    \end{center}
  }{SoftBorder}{SoftWhite}{0.98\linewidth}

  \vspace{0.04in}

  \panel{
    {\large\bfseries\color{TitleBlue} 처리 흐름}

    \vspace{0.02in}
    {\small
    \(
      \texttt{MuJoCo velocimeter}
      \rightarrow
      \texttt{body frame transform}
      \rightarrow
      \texttt{subtract }\omega \times r
      \rightarrow
      \texttt{ROS2 / SITL output}
    \)

    \vspace{0.02in}
    \textbf{코드 포인트}
    \begin{itemize}[leftmargin=1.25em,itemsep=0.18em,topsep=0.22em]
      \item \texttt{ros2\_bridge.py}: \texttt{\_dvl\_velocity\_body()}
      \item \texttt{ros2\_bridge.py}: \texttt{vel\_body = vel\_body - np.cross(gyro\_bmj, r\_body)}
      \item \texttt{tank\_current\_scene.xml}: \texttt{velocimeter dvl\_vel\_body}
    \end{itemize}

    \vspace{0.02in}
    \textbf{요약}: 회전 때문에 생긴 센서 위치 속도를 빼야 DVL 병진 속도와 odometry drift가 더 정확해진다.
    }
  }{TitleBlue}{SoftBlue}{0.98\linewidth}
\end{minipage}

\end{document}
